% Imported from LaTeX_and_pdfs/DGDT_dump.tex; see provenance.json.
\part{Pr\'{a}staro}


Pr\'{a}staro (1996) \cite{dgdt:Pras1996}

\subsubsection{Affine Spaces}

cf. Sec. 1.2 - \emph{Affine Spaces} of Pr\'{a}staro (1996) \cite{dgdt:Pras1996}

\begin{definition}[affine space]
  \begin{equation}
\begin{gathered}
    \text{ affine space \qquad \, } (M, \mathbf{M}, \alpha )  \\
    \text{ with } \\
    \begin{aligned}
      & M \equiv \text{ set (set of pts.) }  \\ 
      & \mathbf{M} \equiv \text{ vector space (space of free vectors) } \\
      & \alpha \equiv \mathbf{M} \times M \to M \equiv \text{ translation operator } \\
      & \alpha : (v,p ) \mapsto p' \equiv p + v
      \end{aligned}
\end{gathered}
  \end{equation}
  Note: $\alpha$ is a \textbf{transitive} action and without fixed pts. (free).

  i.e. $\forall \, p \in M$, 
  \end{definition}

$\forall \, $ pt. $O \in M$, $\alpha:(v,O) \mapsto O' \equiv O + v$, $\alpha (\cdot , O) \equiv \alpha_O \equiv \alpha(O)$.  $\alpha_O(v) = O' = O + \mathbf{v}$ \qquad \, $\forall \, O' \in M$, $\exists \, \mathbf{v} \in \mathbf{M}$ s.t. $O' = O + \mathbf{v}$ \\
$\Longrightarrow M \equiv \mathbf{M}$.

$\forall \, (O, \lbrace e_i \rbrace)_{1 \leq i \leq n }$, where $\lbrace e_i \rbrace$ basis of $\mathbf{M}$, $M \equiv \mathbf{M} = \mathbb{R}^n$ so isomorphism $M \simeq \mathbb{R}^n$ \\

i.e. $\alpha$ is \textbf{without fixed pts.}, meaning, 

Given pointed space $(M, O)$, where base pt. $O\in M$, we can associate $\forall \, p \in M$, vector $\mathbf{x} \in \mathbf{M}$, by 1-to-1 mapping $M\to \mathbf{M}$. 

So for 
\[
\begin{aligned}
& \alpha : \mathbf{M} \times M \to M \\ 
& \alpha(\mathbf{x}, p) = p' = p + \mathbf{x}
\end{aligned}
\]
Consider 
\[
\begin{gathered}
	\alpha(\mathbf{x}, O) = p = \alpha_O(\mathbf{x}) = p \Longrightarrow \exists \, \alpha_O^{-1}(p) = \mathbf{x} \in \mathbf{M}
\end{gathered}
\]

\begin{enumerate}
	\item tangent space of $M$ in $p\in M$ is vector space $T_pM \equiv (\mathbf{M}, p) \cong M$
	\item If $\mathbf{M}$ Euclidean space, affine space $(M, \mathbf{M}, \alpha)$ is Euclidean
	\item Call dim. of affine space $(M, \mathbf{M}, \alpha)$, dim. of $\mathbf{M} \equiv \text{dim}{\mathbf{M}}$
\end{enumerate}
$\lbrace \mathbf{e}_i \rbrace$ basis of $\mathbf{M}$




\begin{definition}
  $(O, \lbrace e_i \rbrace) \equiv $ affine frame.

  $\forall \, $ affine frame $(O,\lbrace e_i \rbrace)$, $\exists \, $ coordinate system $x^{\alpha} : M \to \mathbb{R}$, \\
  where $x^{\alpha}(p)$ is $\alpha$th component, in basis $\lbrace e_i \rbrace$, of vector $p-O$
  \end{definition}

\begin{proposition}[1.6, Pr\'{a}staro (1996) \cite{dgdt:Pras1996}]
	$\forall \, O \in M$, we have canonical identification $M \equiv \mathbf{M}$, since
	\[
	\begin{gathered}
	\begin{aligned}
	& \alpha^{-1}_O:M \to \mathbf{M} \\
	& \alpha^{-1}_O(p) = \mathbf{x} 
	\end{aligned} \qquad \quad \, 
	\begin{aligned}
	& \alpha_O : \mathbf{M} \to M \\
	& \alpha_O: \mathbf{x} = \alpha(\mathbf{x}, O) = p
	\end{aligned}
	\end{gathered}
	\]
	Furthermore, \\
	$\forall \, $ \textbf{affine frame} $(O, \lbrace \mathbf{e}_i \rbrace)_{1\leq i \leq d}$, where $\lbrace \mathbf{e}_i \rbrace$ basis of $\mathbf{M}$, \\
	\phantom{Furthermore} $\exists \, $ isomorphism $M \cong \mathbb{R}^d$, \\
	Then, $\forall \, (O, \lbrace \mathbf{e}_i \rbrace)_{1 \leq i \leq d}$, \\
	\phantom{Then} $\exists \, $ coordinate system $x^{\alpha} : M \to \mathbb{R}$, \\
	\phantom{Then $\exists \, $} where $x^{\alpha}(p) = \alpha$th component, in basis $\lbrace \mathbf{e}_i \rbrace$, of vector $p - O$. 
	
	
	
\end{proposition}

\begin{theorem}[1.4 Pr\'{a}staro (1996) \cite{dgdt:Pras1996}]
  Let $(x^{\alpha}), (\overline{a}^{\alpha})$ 2 coordinate systems correspond to affine frames $(O, \lbrace e_i \rbrace)$, $( \overline{O}, \lbrace \overline{e}_i \rbrace )$, respectively.
\begin{equation}
  \overline{x}^{\alpha} = A^{\alpha}_{ \, \, \beta} x^{\beta} + y^{\alpha}
\end{equation}
where
\[
y^{\alpha} \in \mathbb{R}^n, \qquad \, A^{\alpha}_{ \, \, \beta} \in GL(n; \mathbb{R})
\]
\end{theorem}

\begin{definition}[1.10 Pr\'{a}staro (1996) \cite{dgdt:Pras1996}]
  \begin{equation}
    A(n) \equiv Gl(n,\mathbb{R}) \times \mathbb{R}^n
  \end{equation}
  affine group of dim. $n$
  \end{definition}

\begin{theorem}[1.5] symmetry group of $n$-dim. affine space, called affine group $A(M)$ of $M$.  $\exists \, $ isomorphism,
  \begin{equation}
    A(M) \simeq A(n), \qquad \, f\mapsto (f^{\alpha}_{ \, \, \beta} , y^{\alpha}) \, ; \qquad \, f^{\alpha} \equiv x^{\alpha} \circ f = f^{\alpha}_{ \, \, \beta} x^{\beta} + y^{\alpha}
  \end{equation}
cf. Eq. 1.4 Pr\'{a}staro (1996) \cite{dgdt:Pras1996}
  \end{theorem}

\begin{definition}[metric]
	Let smooth manifold $M$, $\text{dim}M = n$, $\forall \, p \in M$, $\exists \, $ vector space $T_pM$, and so for 
	\begin{equation}
	\begin{aligned}
	& g_p (T_pM)^2 \to \mathbb{R} \\
	& g_p : (X_p,Y_p) \mapsto g_p(X_p,Y_p) \in \mathbb{R}
	\end{aligned}
	\end{equation}
	with $g_p$ being bilinear, symmetric (in $X_p,Y_p$), nondegenerate (i.e. if $g_p(X_p,Y_p)=0$, then $X_p$ or $Y_p=0$)
	
	Note that 
	\[
	g\in \Gamma((TM \otimes TM)^*)
	\]
	and that for $\begin{aligned} & \quad \\ 
	& X = X^i \frac{ \partial }{ \partial x^i } \\
	& Y = Y^i \frac{ \partial }{ \partial x^i } \end{aligned}$  
	so 	
	\[
	g(X,Y) = g_{ij} X^i Y^j
	\]  
	
\end{definition}

Now for 
\[
\begin{gathered}
\begin{aligned} 
& F: M \to N \\
& F: x \mapsto y = y(x) 
\end{aligned} \qquad \, \begin{aligned}
& DF \equiv F_* : T_p M \to T_{F(p)}N \\ 
&  DF : X_p \mapsto (DF) (X^j \frac{ \partial }{ \partial x^j }) = X^j \frac{\partial y^i}{\partial x^j} \frac{ \partial }{ \partial y^i}
\end{aligned}
\end{gathered}
\]
\[
\begin{gathered}
(F^* g')(X,Y) = (F^*g')(X^i \frac{ \partial }{ \partial x^i } , Y^j \frac{ \partial }{ \partial x^j} ) = (F^*g')_{ij} X^i Y^j = g'(F_*X, F_*Y) = \\
= 
\end{gathered}
\]

